\documentclass[11pt]{article}
\usepackage[a4paper,margin=25mm]{geometry}
\usepackage{fontspec}
\usepackage[hidelinks]{hyperref}
\setmainfont{Latin Modern Roman}
\pagestyle{plain}
\begin{document}
{\LARGE\bfseries ALIUS Bulletin - Guidelines for podcast - 2022\par}
\vspace{1em}
\section*{Alius Bulletin}
\section*{*}
\section*{Guidelines For Podcast}
\section*{“What Is It Like ?”}
If you want to conduct a “what is it like” podcast, please read the following guidelines carefully.\par

Form of the podcast episode\par

The podcast aims at lasting between 5 and 10 minutes. The author provides a written text (between 500 and 1000 words) with bibliographic sources that is turned into an audio format by editors. It is read by an English speaker with an audio soundtrack added in the background. Timestamps are used for quoting the scientific literature.\par

Content of the podcast episode\par

“What is it like to be…?” intends to provide an overview of the scientific uptake about what it is like to be in a certain conscious state. It thus favors an explanation of the phenomenological features of this conscious state over mechanistic explanations or functional accounts. It needs to be anchored into the scientific literature but remains accessible to a large audience and avoids unexplained technical terms. Personal and anecdotal takes can be provided as long as it cannot be confounded with scientific data. Historical insights about how the scientific understanding of the phenomenology of this state has evolved over time as well as a discussion of the current research questions and experimental challenges are welcome. 
A full bibliography and quotations in the text are provided by authors. Please find here an example if needed (“what is it like to be asleep”: text and episode).
The podcast is shared on youtube and podcast platforms. A soundtrack is used to give a subjective rendition of what it is like to be in the given conscious state. The soundtrack will be used to give rhythm to the episode and will be tuned down at certain points to highlight important excerpts of the text.\par

Modus operandi and requirements\par

1. Inform the podcast editors that you want to write about what it is like to be in a conscious state.\par

2. Send to the editors your text (500-1000 words) with quotations and full bibliography.\par

3. Editors proofread your text and suggest edits to improve understanding and flow.\par

4. Provide a revised edition of the text based on editors’ suggestions.\par

5. Once the text is validated, a speaker is agreed upon by writers and editors and the audio version of the podcast is recorded. No further changes to the text will be possible at this point. The podcast is recorded, the soundtrack is created to fit the audio and a thumbnail is added.\par

6. The final version of the podcast is sent to authors for approval.\par

7. The podcast episode is published depending on the editorial agenda.

For any questions, feel free to ask editors (here below)\par

Editorial team
Léna Coutrot, George Fejer, Daniel Friedman, Matthieu Koroma\par

\end{document}
